\begin{textsample}{POS Dim 2 – ms – Score 22.00 – ms/ms\_em\_34.txt}  \label{ex:f2_pos_044}
4.1. \textbf{Itinerários} \textbf{Formativos} As mudanças ocorridas na sociedade ao longo dos anos, resultantes do avanço tecnológico e das relações estabelecidas em uma sociedade globalizada, ocasionam impactos nos mais diversos segmentos, inclusive na educação.
Assim, o desafio está em ressignificar a educação pública brasileira, em especial a etapa do Ensino Médio, por meio de um ensino que proporcione uma aprendizagem direcionada para a formação integral do estudante.
A necessidade de estruturas e percursos curriculares dotados de flexibilidade, a apropriação de recursos pedagógicos tecnologicamente avançados e as novas demandas de mercado, somadas a tantos outros fatores, constituem um desafio para qualquer sociedade, particularmente para as instituições associadas à educação.
No enfrentamento desse desafio, o Ensino Médio acha -se diante de uma \textbf{missão} de grandes proporções no sentido de possibilitar uma formação pertinente aos novos tempos, de aumentar as taxas de desempenho escolar, de difundir de forma significativa a chamada propensão para aprender e de garantir a relação da escola com o mundo do trabalho.
Associado a tais fatores, o atual cenário educacional constitui -se em momento apropriado e oportuno para tornar o Ensino Médio mais significativo e pertinente aos tempos atuais.
Nesse sentido, a LDB define a estrutura para o \textbf{currículo} do Ensino Médio em Base Nacional Comum Curricular e \textbf{Itinerários} \textbf{Formativos}, e estabelece que os \textbf{Itinerários} \textbf{Formativos} devem ser organizados por meio da \textbf{oferta} de diferentes arranjos curriculares, integrados ou não, e que considerem as Áreas de Conhecimento ou de atuação \textbf{profissional}, a saber : Linguagens e suas Tecnologias, Matemática e suas Tecnologias, Ciências da Natureza e suas Tecnologias, Ciências Humanas e Sociais Aplicadas e a Formação \textbf{Técnica} Profissional.
Esses \textbf{Itinerários}, com suas respectivas habilidades e competências, regulamentadas por meio da Portaria/MEC n. 1.432, de 28 de dezembro de 2018, admitem que as combinações de arranjos curriculares sejam bem amplas.
A flexibilidade e formas diferenciadas na organização curricular permitem que o estudante tenha diferentes perspectivas, opções de escolha e avanços durante sua \textbf{trajetória} acadêmica por meio de ênfases, temas integradores, aprofundamentos, estudos complementares e progressividade dos conteúdos em determinadas áreas do conhecimento ou de forma integrada.
Tal organização torna imprescindível que os sistemas de ensino, \textbf{redes} escolares e escolas reorientem seus \textbf{currículos} e projetos político-pedagógicos, visto que “ a \textbf{carga} \textbf{horária} mínima anual do Ensino Médio deve ser ampliada de forma progressiva para mil e quatrocentas horas, devendo os sistemas de ensino oferecer no prazo máximo de cinco anos, pelo menos mil horas anuais, a partir de dois de março de 2017 ”. ( BRASIL, 1996 ).
Nessa nova organização, a \textbf{carga} \textbf{horária} \textbf{destinada} ao cumprimento da Base Nacional Comum Curricular está definida em até 1.800 ( mil e oitocentas ) horas.
Em uma \textbf{oferta} de 3.000 ( três mil ) horas, 1.800 ( um mil e oitocentas ) horas da \textbf{carga} \textbf{horária} são \textbf{destinadas} à Formação Geral Básica e 1.200 ( mil e duzentas ) horas, para os \textbf{Itinerários} \textbf{Formativos}.
A Resolução CNE/CEB n. 3, de 21 de novembro de 2018, estabelece em seu artigo 12 : > § 6º Os sistemas de ensino devem garantir a \textbf{oferta} de mais de um \textbf{itinerário} \textbf{formativo} em cada \textbf{município}, em áreas distintas, permitindo -lhes a escolha, dentre diferentes arranjos curriculares, \textbf{atendendo} assim a heterogeneidade e pluralidade de condições, interesses e aspirações ( BRASIL, 2018a ).
Essa afirmação vem apoiar o princípio da flexibilização curricular para os estudantes da etapa do Ensino Médio e considera, sobretudo, as particularidades e os arranjos produtivos de cada região/município, no intuito de garantir a possibilidade de escolha de diferentes \textbf{Itinerários} \textbf{Formativos}.
A flexibilidade em questão implica na organização do processo ensino-aprendizagem, de modo a permitir que o estudante tenha diferentes perspectivas durante a sua \textbf{trajetória} escolar. \textbf{Itinerários} \textbf{Formativos} \textbf{flexíveis} objetivam o aprofundamento de estudos, o avanço ou a \textbf{garantia} de estudos de complementação em determinadas Áreas de Conhecimento, com vistas a promover uma aprendizagem com maior profundidade e que estimule o desenvolvimento integral do estudante, por meio do incentivo ao protagonismo, à autonomia e à responsabilidade do estudante por suas escolhas e seu futuro.
Para orientar os sistemas de ensino na construção dos \textbf{Itinerários} \textbf{Formativos}, o Ministério da Educação ( MEC ) publicou, em 28 de dezembro de 2018, por meio da \textbf{Portaria} n. 1.432, os Referenciais Curriculares para a Elaboração de \textbf{Itinerários} \textbf{Formativos}, documento em consonância com a Resolução CNE/CEB n. 3, de 21 de novembro de 2018, que estabelece em seu artigo 7{\EmojiFont °} : > § 2º os \textbf{itinerários} \textbf{formativos} orientados para o aprofundamento e ampliação das aprendizagens em áreas do conhecimento devem garantir a apropriação de procedimentos cognitivos e uso de metodologias que favoreçam o protagonismo juvenil, e organizar -se em torno de um ou mais dos seguintes eixos \textbf{estruturantes} : > > - I.
Investigação científica : supõe o aprofundamento de conceitos fundantes das ciências para a interpretação de ideias, fenômenos e processos para serem utilizados em procedimentos de investigação voltados ao enfrentamento de situações cotidianas e demandas locais e coletivas, e a proposição de intervenções que considerem o desenvolvimento local e a melhoria da \textbf{qualidade} de vida da comunidade ; > - II.
Processos criativos : supõem o uso e o aprofundamento do conhecimento científico na construção e criação de experimentos, modelos, protótipos para a criação de processos ou produtos que \textbf{atendam} a demandas pela resolução de problemas identificados na sociedade ; > - III.
Mediação e intervenção sociocultural : supõem a mobilização de conhecimentos de uma ou mais áreas para medir conflitos, promover entendimento e implementar soluções para questões e problemas identificados na comunidade ; > - IV. \textbf{Empreendedorismo} : supõe a mobilização de conhecimentos de diferentes áreas para a formação de organizações com variadas \textbf{missões} voltadas ao desenvolvimento de produtos ou prestação de serviços inovadores com o uso das tecnologias ( BRASIL, 2018a ).
Evidencia -se, assim, que a construção dos \textbf{Itinerários} \textbf{Formativos} passe, necessariamente, por um eixo \textbf{estruturante} ou, preferencialmente, por todos, uma vez que desenvolvem habilidades importantes para a formação integral dos estudantes e conectam experiências educativas com a realidade contemporânea, por meio de diferentes arranjos curriculares.
A reorganização do Ensino Médio deve promover a difusão do conhecimento por meio da Formação Geral Básica e de \textbf{Itinerários} \textbf{Formativos}, incorporar questões de \textbf{empreendedorismo} e de inovação e garantir um \textbf{currículo} que viabilize o diálogo entre a educação e o mundo do trabalho.
O desafio está em \textbf{flexibilizar} os \textbf{currículos}, de modo a substituir a noção de \textbf{curso} por percurso e, ao mesmo tempo, enfatizar a formação cidadã e a formação de valores, reconhecer que a escola deve incentivar o protagonismo estudantil e a sua coautoria no processo de aprendizagem, assim como proporcionar a articulação de estratégias e instrumentos que permitam novas atuações no campo da cidadania.

% matched lemmas: atender, carga, currículo, curso, destinar, empreendedorismo, estruturante, flexibilizar, flexível, formativo, garantia, horário, itinerário, missão, município, oferta, portaria, profissional, qualidade, rede, trajetória, técnico
\end{textsample}
